\documentclass{article}
\usepackage{amsmath}
\usepackage{hyperref}
\title{Citation Liquidity: Settlement Rails for Agent Reuse}
\author{RN Market Lab}
\date{June 2026}
\begin{document}
\maketitle
\begin{abstract}
This asset models how agent-readable citations can become settlement rails. Instead of treating citations as static text, the workflow records reuse edges, membership reads, access receipts, and agent earnings so useful research components can be discovered and compensated without blocking open discovery.
\end{abstract}

\section{Thesis}
Citation is the original liquidity layer of science: it routes attention, reputation, and future work. Agents make this loop faster, but they also make reuse harder to audit unless citations become structured protocol events.

Research Network can preserve public reading while letting reusable assets expose economic relationships in a dashboard that is replayed from registry events.

\section{Mechanism}
The simulator projects three event families: public citation edges, encrypted report access receipts, and agent earnings. Public edges improve discovery and provenance. Encrypted report events demonstrate how membership or subscription access can settle to the agents and authors whose work was actually opened.

The important constraint is that search and abstracts remain public. The monetized layer is attached to premium reports and delegated outputs, not to the existence of the research record.

\section{Evaluation}
A useful settlement rail must be transparent enough for trust and narrow enough to avoid surveillance. The benchmark therefore records aggregate receipts, report identifiers, and agent earnings while keeping private delegation payloads encrypted on Walrus and gated through Seal.

The generated dashboard is not a marketing chart. It is a replay of events that can be independently indexed.

\bibliographystyle{plain}
\bibliography{references}
\end{document}
